\section{Discussion}
\label{sec:discussion}

\subsection{CPA vs.\ DPA: Efficiency}

Both attacks recover the full 128-bit AES key from 2{,}000 traces on the
reference dataset.
However, the two attacks differ in the structure of their sweep:

\begin{itemize}
  \item \textbf{CPA} requires a sweep of the leakage constant $c$
        (0--255) to handle the HD leakage of byte 11.
        The HW model ($c=0$) already gives rank~1 for 15/16 bytes,
        and $c=2$ resolves the remaining byte.
        The Pearson correlation distinguisher is strictly more data-efficient
        than DPA because it uses the \emph{magnitude} of the hypothetical
        power (not just a single bit), exploiting more signal per trace.

  \item \textbf{DPA} requires a bit-position sweep ($\ell = 0$--7).
        Bit 1 of the last-round intermediate produces clean partitioning for
        all 16 bytes simultaneously.
        DPA is conceptually simpler and makes no assumption about the leakage
        model beyond a single bit's influence on power.
\end{itemize}

\subsection{Byte 11: Hamming Distance Leakage}

The most notable finding is that byte 11 of the round-10 key does not follow
a pure HW leakage model.
Under $c=0$, the CPA score places the correct hypothesis at rank 2
rather than rank 1; $c=2$ (equivalent to $\mathrm{HW}(v \oplus 2)$) resolves
this.
This indicates that the power consumed during the byte-11 operation is
proportional to the number of bit transitions \emph{from state 2} to the
computed intermediate value---a typical HD leakage pattern caused by register
pipeline state.
DPA is unaffected by this because it uses a single bit of $v$ as the
distinguisher, not its full value.

\subsection{Relationship to the AES S-Box Structure}

The attack targets the SubBytes operation because the S-Box is the only
non-linear component in AES.
This is where plaintext (or ciphertext) and key bytes are combined in a way
that is both \emph{computable from known data} (once a key hypothesis is
fixed) and \emph{physically observable} (the S-Box lookup causes data-
dependent switching).
Linear operations (ShiftRows, MixColumns, AddRoundKey) still leak, but their
leakage is linearly predictable from the SubBytes output and thus not
independently useful.

\subsection{Limitations and Open Questions}

\begin{enumerate}
  \item \textbf{Real hardware GE curves.}
    The reference dataset traces have $T = 29{,}000$ samples; the Pico
    acquisition is configured for $T = 5{,}000$ samples.
    The GE analysis (Section~\ref{sec:results:ge}) on real Pico traces
    is the primary pending deliverable.

  \item \textbf{Masking countermeasure.}
    The current firmware has no countermeasures.
    Applying first-order Boolean masking to the S-Box lookup would break
    first-order CPA/DPA and requires a second-order attack.
    This is a natural extension.

  \item \textbf{Alignment sensitivity.}
    Traces must be aligned on the trigger edge for both attacks to function.
    Even a one-sample misalignment spreads the leakage signal across adjacent
    samples and reduces distinguisher scores.
    The hardware trigger (GPIO\,22) provides sample-accurate alignment for
    this setup.
\end{enumerate}
